\section{Stable unlabelled unicyclic graphs: \texorpdfstring{\seq{A006545}}{A006545}}
\label{sec:a006545}

A graph is semistable at a vertex \(v\) when every automorphism of \(G-v\)
preserves the former neighborhood of \(v\), and is stable when its vertices
admit a deletion ordering in which the current graph is semistable at the next
vertex.  McAvaney, Grant, and Holton proved that a connected unicyclic graph is
stable if and only if its automorphism group contains a transposition
\citep{McAvaneyGrantHolton1974,McAvaney1975}.  This characterization turns the
enumeration into a local classification problem.

\subsection{Rooted-tree series and decorated cycles}

Let \(R(x)\) be the ordinary generating function for unlabelled rooted trees.
The classical Pólya-Otter equation is
\begin{equation}
 R(x)=x\exp\left(\sum_{j\ge1}\frac{R(x^j)}{j}\right).
 \label{eq:a006545-R}
\end{equation}
Let \(B(x)\) count rooted trees in which no vertex has two or more leaf
children.  In the rooted-tree multiset construction, arbitrary multiplicity of
leaf children contributes \((1-x)^{-1}\), while multiplicity zero or one
contributes \(1+x\).  Replacing the former by the latter multiplies the rooted
operator by \((1+x)(1-x)=1-x^2\), giving
\begin{equation}
 B(x)=x(1-x^2)\exp\left(\sum_{j\ge1}\frac{B(x^j)}{j}\right).
 \label{eq:a006545-B}
\end{equation}

For a series \(F\) with \(F(0)=0\), let \(D(F)\) be the OGF for unoriented
cycles of length at least three whose vertices carry \(F\)-objects.  Burnside's
lemma for the dihedral group gives
\begin{align}
 D(F)&=\sum_{k\ge3}\frac{\operatorname{Rot}_k(F)+
                              \operatorname{Ref}_k(F)}{2k},
 \label{eq:a006545-D}\\
 \operatorname{Rot}_k(F)&=\sum_{d\mid k}\varphi(d)F(x^d)^{k/d},
 \label{eq:a006545-rot}\\
 \operatorname{Ref}_{2m+1}(F)&=(2m+1)F(x)F(x^2)^m,
 \label{eq:a006545-refodd}\\
 \operatorname{Ref}_{2m}(F)&=m\left(F(x)^2F(x^2)^{m-1}+F(x^2)^m\right).
 \label{eq:a006545-refeven}
\end{align}
Every connected unicyclic graph has a unique cycle, with a rooted tree attached
at each cycle vertex.  Consequently \(D(R)\) counts all connected unlabelled
unicyclic graphs.

\subsection{Where a transposition can occur}

\begin{lemma}[Twin classification]
\label{lem:a006545-twins}
If a connected unicyclic graph has an automorphism that transposes exactly two
vertices \(u,v\), then the pair has exactly one of the following forms.
\begin{enumerate}
\item \(u,v\) are two leaf children of the same vertex;
\item \(u,v\) are the two bare vertices of a triangle;
\item \(u,v\) are opposite bare vertices of a 4-cycle.
\end{enumerate}
Conversely, each of these configurations yields a transposition.
\end{lemma}

\begin{proof}
A transposition fixing every other vertex implies
\[
 N(u)\setminus\{v\}=N(v)\setminus\{u\};
\]
thus \(u,v\) are twins.  Suppose first that they are nonadjacent.  If they have
one common neighbor, any additional incident edge would also be common and
would create a second cycle, so they are sibling leaves.  If they have two
common neighbors, the four vertices form the unique cycle, which is a 4-cycle,
and the transposed vertices are bare.  Three common neighbors create at least
two cycles.

If \(u,v\) are adjacent, connectedness of a nontrivial unicyclic graph forces a
common neighbor.  There can be exactly one: no common neighbor would isolate a
\(K_2\) component, while two would create more than one cycle.  The pair is
therefore the two bare vertices of a triangle.  In all three cases the indicated
swap visibly preserves adjacency and fixes every other vertex.
\end{proof}

\subsection{The ordinary generating function}

The difference \(D(R)-D(B)\) counts exactly the unicyclic graphs containing a
pair of sibling leaves.  Among the remaining \(B\)-decorated cycles, a triangle
transposition has two bare vertices and an arbitrary \(B\)-object at the third,
contributing \(x^2B(x)\).  A 4-cycle transposition has two opposite bare
vertices; the other opposite pair carries an unordered pair of \(B\)-objects,
whose OGF is
\[
 \frac{B(x)^2+B(x^2)}{2}.
\]
The triangle and 4-cycle families are disjoint and contain no sibling-leaf
transposition.

\begin{theorem}[All-\(n\) OGF]
\label{thm:a006545}
The ordinary generating function for stable connected unlabelled unicyclic
graphs is
\begin{equation}
 \boxed{
 A(x)=D(R)-D(B)+x^2B(x)+
       \frac{x^2}{2}\left(B(x)^2+B(x^2)\right),}
 \label{eq:a006545-main}
\end{equation}
where \(R,B,D\) are determined by
\cref{eq:a006545-R,eq:a006545-B,eq:a006545-D,eq:a006545-rot,eq:a006545-refodd,eq:a006545-refeven}.
Thus \(a(n)=[x^n]A(x)\) for every \(n\ge3\).
\end{theorem}

\begin{proof}
By the cited structural theorem, stability is equivalent to the existence of a
transposition.  Lemma~\ref{lem:a006545-twins} partitions such graphs into those
with sibling leaves and the two exceptional bare-cycle mechanisms.  The four
terms in \cref{eq:a006545-main} count these disjoint classes exactly.
\end{proof}

The coefficients begin, with offset three,
\[
\begin{aligned}
1,&\ 2,\ 3,\ 8,\ 22,\ 62,\ 176,\ 500,\ 1425,\ 4078,\\
&11666,\ 33447,\ 95922,\ 275332,\ 790518,\ 2270963,\\
&6525674,\ 18759532,\ldots.
\end{aligned}
\]
The exact series code independently reproduces the total unicyclic counts
through 12 vertices.  Complete graph6 enumerations of connected unlabelled
\(n\)-vertex, \(n\)-edge graphs were checked in three ways: by the recursive
stability definition, by direct automorphism-group sizes at every deletion
state, and by the elementary transposition test of
\cref{lem:a006545-twins}.  All three select the same records through
\(n=12\).
